% ============================================================
%  Part 16 — Adiabatic Cooling, Photon Gas, Grand Canonical
%  Transcribed from pages 201–215
% ============================================================

% ------------------------------------------------------------------
%  Continuation: high-T limit, heat capacity, and adiabatic cooling
% ------------------------------------------------------------------

\subsection*{High-temperature limit of the Einstein solid}

In the limit $T \to \infty$ we have $x \equiv \beta\hbar\omega \ll 1$, so
\begin{equation}
    \frac{1}{e^x - 1} \approx \frac{1}{x + \tfrac{1}{2}x^2}
    = \frac{1}{x\!\left(1+\tfrac{1}{2}x\right)} \approx \frac{1}{x}\!\left(1-\frac{x}{2}\right)
    \approx \frac{1}{x},
\end{equation}
and therefore the internal energy becomes
\begin{equation}
    U(T\to\infty) = \frac{3N\hbar\omega}{\beta\hbar\omega} = 3Nk_BT.
\end{equation}
This agrees with the equipartition theorem: all $6$ degrees of freedom (three kinetic, three potential) of the oscillators are active, each contributing $\tfrac{1}{2}k_BT$.

\subsection*{Heat capacity of the Einstein solid}

Differentiating the internal energy with respect to temperature yields
\begin{equation}
    C_V = \pdv{U}{T} = 3Nk_B\!\left(\frac{\hbar\omega}{k_BT}\right)^{\!2}
    \frac{e^{\beta\hbar\omega}}{\left(e^{\beta\hbar\omega}-1\right)^2}.
\end{equation}

\paragraph{Low-temperature limit ($T\to 0$).}  As $T\to 0$ the exponential
$e^{\beta\hbar\omega}$ dominates the denominator, giving
\begin{equation}
    C_V \approx 3Nk_B\!\left(\frac{\hbar\omega}{k_BT}\right)^{\!2}
    e^{-\hbar\omega/k_BT} \xrightarrow{T\to 0} 0.
\end{equation}
All degrees of freedom are frozen out: the thermal energy $k_BT$ is far too small
to excite any oscillator across the gap $\hbar\omega$.

\paragraph{High-temperature limit ($T\to\infty$).}  Expanding the denominator for small
$\beta\hbar\omega$ recovers
\begin{equation}
    C_V \to 3Nk_B.
\end{equation}

\begin{fact}
\textbf{Dulong--Petit Law.}  At high temperatures the heat capacity of a monatomic
solid saturates at $C_V = 3Nk_B = 3R$ per mole, corresponding to the full
equipartition result for three-dimensional harmonic oscillators.
\end{fact}

% ------------------------------------------------------------------
\section{Final Example: Adiabatic Cooling}
% ------------------------------------------------------------------

We consider $N$ classical, indistinguishable, relativistic particles confined in a
linear (absolute-value) potential.  Specifically:
\begin{itemize}
    \item \textit{Classical:} the typical number of particles per quantum state is
          much less than one.
    \item \textit{Indistinguishable:} the $N$-particle partition function acquires
          the Gibbs correction $Z_N = Z_1^N / N!$.
    \item \textit{Relativistic with linear potential:} single-particle energy is
          $E = c|p|$ and potential energy is $V(x)=\alpha|x|$.
\end{itemize}

The single-particle partition function is
\begin{equation}
    Z_1 = \int \frac{d^3x\, d^3p}{(2\pi\hbar)^3}\;
    e^{-\beta c|p|}\, e^{-\beta\alpha|x|}.
\end{equation}

\paragraph{Evaluating the integrals.}  Both the momentum and position integrals are
of the form $\int d^3u\, e^{-a|u|}$.  In spherical coordinates,
\begin{align}
    \int d^3u\; e^{-a|u|}
    &= 4\pi \int_0^\infty u^2 e^{-au}\,du
     = 4\pi a^{-3} \int_0^\infty r^2 e^{-r}\,dr
     = \frac{8\pi}{a^3},
\end{align}
using $\int_0^\infty r^2 e^{-r}\,dr = 2$.  Applying this to both factors gives
\begin{equation}
    Z_1 = \frac{1}{(2\pi\hbar)^3}(8\pi)^2
    \left(\frac{k_BT}{c}\right)^{\!3}\!\left(\frac{k_BT}{\alpha}\right)^{\!3}
    = \frac{8}{\pi}\!\left(\frac{1}{\beta c\alpha}\right)^{\!3}(k_BT)^6.
\end{equation}

The internal energy and entropy of the $N$-particle system follow from
$U_N = -\partial_\beta \log Z_N$ and $S_N = U_N/T + k_B\log Z_N$:
\begin{align}
    U_N &= 6Nk_BT,\\
    S_N &= 6Nk_B + k_B\!\left[N\log Z_1 - N\log N + N\right].
\end{align}
The factor $6$ confirms that each particle has six effective degrees of freedom
(three from the relativistic momentum integral, three from the linear position
potential).

\paragraph{Adiabatic change of the trap gradient.}  If $\alpha$ is changed
\emph{adiabatically}, then $S_N$ is held constant.  Since $S_N$ depends on $T$
and $\alpha$ only through $Z_1 \propto T^6/\alpha^3$, holding $S_N$ constant
requires $Z_1 = \text{const}$, i.e.\
\begin{equation}
    \frac{T^6}{\alpha^3} = \text{const} \implies T^2 \propto \alpha
    \implies \boxed{T \sim \sqrt{\alpha}}.
\end{equation}
Reducing the trap gradient adiabatically therefore cools the gas: as $\alpha$ is
slowly decreased, the temperature drops as $\sqrt{\alpha}$.

% ------------------------------------------------------------------
\section{Electromagnetic Radiation: the Photon Gas}
% ------------------------------------------------------------------

\subsection{Physical setup}

Consider an empty cubic box of side $L$ held at temperature $T$.  The walls
continuously absorb and emit photons, maintaining thermal equilibrium.  We want
to find the allowed modes of the electromagnetic field inside, quantise them, and
compute the thermodynamics.

\subsection{Allowed modes from Maxwell's equations}

In vacuum with no charges or currents, Maxwell's equations give
\begin{align}
    \left(\frac{1}{c^2}\partial_t^2 - \nabla^2\right)\vec{E} &= 0, &
    \nabla\cdot\vec{E} &= 0,\\
    \nabla\cdot\vec{B} &= 0, &
    \frac{1}{c}\partial_t\vec{B} + \nabla\times\vec{E} &= 0\quad\text{(Faraday)}.
\end{align}

Seeking standing-wave solutions satisfying the conducting boundary condition
$E_\parallel = 0$ at each wall, the components of $\vec{E}$ take the form
\begin{align}
    E_x &= E_{x0}\cos(k_x x)\sin(k_y y)\sin(k_z z)\sin(\omega t),\notag\\
    E_y &= E_{y0}\sin(k_x x)\cos(k_y y)\sin(k_z z)\sin(\omega t),\\
    E_z &= E_{z0}\sin(k_x x)\sin(k_y y)\cos(k_z z)\sin(\omega t).\notag
\end{align}
The dispersion relation is $\omega = c|\vec{k}|$.  The transversality condition
$\nabla\cdot\vec{E}=0$ requires $k_xE_{x0}+k_yE_{y0}+k_zE_{z0}=0$, which removes
one component of $\vec{E}_0$, leaving two independent polarisations.

\subsection{Boundary conditions and allowed wavevectors}

Applying the boundary conditions at each pair of walls:
\begin{itemize}
    \item At $x=L$: require $E_y=E_z=0$, so $\sin(k_xL)=0$ giving
          $k_x = n_x\pi/L$.
    \item At $y=L$: $k_y = n_y\pi/L$.
    \item At $z=L$: $k_z = n_z\pi/L$.
\end{itemize}

\begin{formula}
The allowed wavevectors are
\begin{equation}
    \vec{k} = \frac{\pi}{L}(n_x,\,n_y,\,n_z),\qquad n_x,n_y,n_z = 0,1,2,\ldots
\end{equation}
For each $\vec{k}$ there are \emph{two} independent polarisation modes.
\end{formula}

\subsection{Quantum description of photons}

Each EM mode of wavevector $\vec{k}$ satisfies
\begin{equation}
    \partial_t^2 \vec{E} = -\omega^2 \vec{E},
\end{equation}
which is identical to the harmonic oscillator equation.  Quantising each mode,
the energy levels are
\begin{equation}
    \mathcal{E}_n(\vec{k}) = \left(n+\tfrac{1}{2}\right)\hbar\omega(\vec{k}),
    \qquad n = 0,1,2,\ldots,
\end{equation}
where $n$ is the number of photons (quanta of energy) in that mode.

The full state of the EM field is specified by the set of occupation numbers
$\{n_{\vec{k},\lambda}\}$ for every mode $(\vec{k},\lambda)$, where $\lambda=1,2$
labels the two polarisations.

\begin{definition}
The \textbf{total energy of the EM field} in a given state is
\begin{equation}
    E(\text{state}) = \sum_{\vec{k},\lambda}
    \left(n_{\vec{k},\lambda}+\tfrac{1}{2}\right)\hbar\omega(\vec{k}).
\end{equation}
\end{definition}

\subsection{Partition function}

Summing over all configurations $\{n_{\vec{k},\lambda}\}$ and using the fact that
the exponential of a sum factors into a product of exponentials,
\begin{align}
    Z &= \sum_{\{n_{\vec{k},\lambda}\}}
    \exp\!\left\{-\beta\sum_{\vec{k},\lambda}
    \left(n_{\vec{k},\lambda}+\tfrac{1}{2}\right)\hbar\omega(\vec{k})\right\}
    = \prod_{\vec{k},\lambda}\sum_{n=0}^{\infty}
    \exp\!\left(-\beta\!\left(n+\tfrac{1}{2}\right)\hbar\omega(\vec{k})\right).
\end{align}
Each factor is a geometric series, giving
\begin{equation}
    Z = \prod_{\vec{k},\lambda} Z_\omega(\vec{k}),\qquad
    Z_\omega = \frac{e^{-\beta\hbar\omega/2}}{1-e^{-\beta\hbar\omega}}.
\end{equation}

% ------------------------------------------------------------------
\section{EM Radiation Thermodynamics}
% ------------------------------------------------------------------

\subsection{Partition function in the continuum limit}

Including the two polarisations, the log-partition function is
\begin{equation}
    \log Z = \sum_{n_x,n_y,n_z} 2\log Z_\omega.
\end{equation}
In the thermodynamic limit, the sum over non-negative integers becomes an integral
over the positive octant of $n$-space.  Switching to spherical coordinates and
changing variable $\omega = \pi c n/L$ (so $d\omega = \pi c\,dn/L$):
\begin{align}
    \log Z
    &= \frac{4\pi}{8}\int_0^\infty n^2\,dn \cdot 2\log Z_\omega\notag\\
    &= \pi\left(\frac{L}{\pi c}\right)^{\!3}\int_0^\infty \omega^2 \log Z_\omega\,d\omega.
\end{align}
Writing $V=L^3$ and substituting $\log Z_\omega = -\tfrac{\beta\hbar\omega}{2}
-\log(1-e^{-\beta\hbar\omega})$, the zero-point term contributes a (divergent)
constant to $F$ that is the same for all temperatures and can be neglected:

\begin{formula}
\textbf{Partition function of the photon gas (dropping zero-point energy):}
\begin{equation}
    \log Z = \frac{V}{\pi^2 c^3}\int_0^\infty d\omega\;\omega^2
    \left[-\log\!\left(1-e^{-\beta\hbar\omega}\right)\right].
\end{equation}
\end{formula}

\subsection{Internal energy}

\begin{equation}
    U = -\pdv{\log Z}{\beta}
    = \frac{V}{\pi^2 c^3}\int_0^\infty d\omega\;
    \underbrace{\omega^2}_{\text{density of modes}}\;
    \underbrace{\frac{\hbar\omega}{e^{\beta\hbar\omega}-1}}_{\text{mean energy per mode}}.
\end{equation}

To evaluate, substitute $\tau = \beta\hbar\omega$:
\begin{equation}
    U = \frac{\hbar V}{\pi^2 c^3}\!\left(\frac{k_BT}{\hbar}\right)^{\!4}
    \int_0^\infty \frac{\tau^3\,d\tau}{e^\tau-1}
    = \frac{\hbar V}{\pi^2 c^3}\!\left(\frac{k_BT}{\hbar}\right)^{\!4} 3!\,\zeta(4),
\end{equation}
using the standard result $\int_0^\infty \tau^\nu(e^\tau-1)^{-1}\,d\tau = \nu!\,\zeta(\nu+1)$.

\begin{fact}
\textbf{Riemann zeta function values:}
\begin{equation}
    \zeta(2) = \frac{\pi^2}{6},\qquad \zeta(3)\approx 1.202,\qquad \zeta(4)=\frac{\pi^4}{90}.
\end{equation}
\end{fact}

Inserting $3!\,\zeta(4) = 6\cdot\pi^4/90 = \pi^4/15$:

\begin{formula}
\textbf{Energy density of the EM field (Stefan's $T^4$ law):}
\begin{equation}
    \frac{U}{V} = \frac{\pi^2 k_B^4}{15\hbar^3 c^3}\,T^4
    = \frac{4\sigma}{c}\,T^4,
    \qquad \sigma \equiv \frac{\pi^2 k_B^4}{60\hbar^3 c^2}.
\end{equation}
Here $\sigma$ is the \textbf{Stefan--Boltzmann constant}.
\end{formula}

\subsection{Free energy}

Integrating by parts in the frequency integral,
\begin{align}
    F &= -k_BT\log Z
    = \frac{k_BTV}{\pi^2 c^3}\int_0^\infty \omega^2\log(1-e^{-\beta\hbar\omega})\,d\omega\notag\\
    &= \frac{k_BTV}{\pi^2 c^3}\!\left[
      \tfrac{1}{3}\omega^3\log(1-e^{-\beta\hbar\omega})\Big|_0^\infty
      + \frac{\beta\hbar}{3}\int_0^\infty \frac{\omega^3\,d\omega}{e^{\beta\hbar\omega}-1}
    \right]
    = -\frac{U}{3}.
\end{align}
Hence $F = -U/3$.

\subsection{Pressure and entropy}

From $dF = -S\,dT - P\,dV$:
\begin{equation}
    P = -\pdv{F}{V}\bigg|_T = \frac{U}{3V} = \frac{1}{3}\cdot\frac{4\sigma T^4}{c}.
\end{equation}
Radiation pressure is one-third the energy density, analogous to a classical ideal
gas of massless particles.

\begin{formula}
\textbf{Entropy of the photon gas:}
\begin{equation}
    S = -\pdv{F}{T}\bigg|_V = \frac{4U}{3T} = \frac{16\sigma T^3 V}{3c}.
\end{equation}
\end{formula}

\subsection{Number of photons}

Dividing the energy integrand by the energy per photon $\hbar\omega$ gives the
number of photons at each frequency.  Integrating:
\begin{equation}
    N = \frac{V}{\pi^2 c^3}\int_0^\infty \frac{\omega^2\,d\omega}{e^{\beta\hbar\omega}-1}
    = \frac{V}{\pi^2 c^3}\!\left(\frac{k_BT}{\hbar}\right)^{\!3} 2\,\zeta(3)
    \approx \frac{1.48\,\sigma T^3 V}{\hbar c}.
\end{equation}
The number of photons scales as $T^3$ and is not a conserved quantity.

\subsection{Heat capacity and fluctuations}

The heat capacity is
\begin{equation}
    C_V = \pdv{U}{T}\bigg|_V = \frac{16\sigma T^3 V}{c}.
\end{equation}

The relative energy fluctuation scales as
\begin{equation}
    \frac{\sigma(U)}{U} \sim \sqrt{\frac{\hbar c}{\sigma T^3 V}} \sim \frac{1}{\sqrt{N}},
\end{equation}
as expected for a system of $N$ nearly independent modes.

% ------------------------------------------------------------------
\subsection{Some Points on the Photon Gas}
% ------------------------------------------------------------------

\subsubsection*{(1) Mechanical derivation of radiation pressure}

Consider plane waves of mode $\vec{k}$ incident on a wall at angle $\theta$ to
the wall normal.

\begin{center}
[DIAGRAM: Plane wave ray incident at angle $\theta$ to a flat conducting wall;
area element $A$ on wall; hemisphere of radius $c\,\delta t$ indicating the
volume of photons that can reach the wall in time $\delta t$]
\end{center}

In a small time $\delta t$, photons in a volume $A\,c\,\delta t\cos\theta$ heading
toward the wall arrive.  Each mode of energy density $U_k$ contributes energy
$A\,c\,\delta t\cos\theta \cdot U_k$ and hence momentum
$A\,\delta t\cos\theta \cdot U_k$ (since photon momentum equals energy$/c$).
Each reflection bounces the photon back, so the impulse per photon is doubled;
only the component normal to the wall contributes to pressure, giving a factor of
$\cos\theta$ again.  The pressure from one mode is therefore proportional to
$2\cos^2\theta\, U_k$.

Integrating over all incoming directions (hemisphere) and averaging:
\begin{align}
    P &= 2\langle\cos^2\theta\rangle\cdot\frac{U}{V}
    = 2\cdot\frac{\int_0^{\pi/2}\cos^2\theta\sin\theta\,d\theta}
                 {\int_0^{\pi/2}\sin\theta\,d\theta}
    \cdot\frac{U}{2V}
    = \frac{1}{3}\frac{U}{V},
\end{align}
confirming the thermodynamic result.

\subsubsection*{(2) Classical limit: the ultraviolet catastrophe}

As $\hbar\to 0$, the Bose--Einstein factor becomes
$e^{\beta\hbar\omega}-1\approx \beta\hbar\omega$, so the energy density
diverges:
\begin{equation}
    U \to \frac{k_BTV}{\pi^2 c^3}\int_0^\infty \omega^2\,d\omega \to \infty.
\end{equation}
This is the \textbf{Rayleigh--Jeans} result: the classical equipartition theorem
assigns $k_BT$ to every mode, and there are infinitely many high-frequency modes.
The resulting divergence is called the \textbf{ultraviolet catastrophe}, and its
resolution by Planck (through the quantisation of the EM field) was a foundational
moment in quantum mechanics.

\subsubsection*{(3) Spectral emission and the Stefan--Boltzmann law}

Imagine making a small hole in the box.  Photons in a hemisphere of radius $r$
within the small volume $dV = r^2\sin\theta\,d\theta\,d\phi\,dr$ that travel
toward area $A$ of the hole escape if their direction points at the hole and they
are within distance $c\,\delta t$.  The energy emitted per unit time, per unit
area, per unit frequency is
\begin{equation}
    \frac{\text{Power}}{A\;d\omega}
    = \frac{\hbar}{4\pi^2 c^2}\cdot\frac{\omega^3}{e^{\beta\hbar\omega}-1}.
\end{equation}

\begin{fact}
\textbf{Recap --- spectral radiance per unit angular frequency:}
\begin{equation}
    \frac{P/A}{\text{unit}\;\omega}
    = \frac{\hbar}{4\pi^2 c^2}\cdot\frac{\omega^3}{e^{\beta\hbar\omega}-1}.
\end{equation}
\end{fact}

Integrating over all frequencies and using the same zeta-function result:
\begin{align}
    \frac{P}{A}
    &= \frac{\hbar}{4\pi^2 c^2}
       \left(\frac{k_BT}{\hbar}\right)^{\!4}
       \int_0^\infty\frac{\tau^3\,d\tau}{e^\tau-1}
    = \frac{\hbar}{4\pi^2 c^2}
       \left(\frac{k_BT}{\hbar}\right)^{\!4} 3!\,\zeta(4)
    = \frac{\pi^2 k_B^4}{60\hbar^3 c^2}\,T^4.
\end{align}

\begin{formula}
\textbf{Stefan--Boltzmann Law:}
\begin{equation}
    \frac{P}{A} = \sigma T^4,\qquad
    \sigma = \frac{\pi^2 k_B^4}{60\hbar^3 c^2}.
\end{equation}
\end{formula}

\paragraph{Example: Earth's equilibrium temperature.}

The Sun has surface temperature $T_0\approx 5772\,\text{K}$ and radius
$R_\odot\approx 6.96\times 10^8\,\text{m}$; the Earth--Sun distance is
$a\approx 1.5\times 10^{11}\,\text{m}$.  The total solar luminosity is
$P_\odot = 4\pi R_\odot^2\,\sigma T_0^4$.  The power intercepted by Earth's
cross-section $\pi R_E^2$ is
\begin{equation}
    P_\text{in} = P_\odot\cdot\frac{\pi R_E^2}{4\pi a^2}.
\end{equation}
Setting this equal to the power radiated by Earth's full surface,
$P_\text{out} = 4\pi R_E^2\,\sigma T_E^4$, gives

\begin{formula}
\textbf{Earth's equilibrium temperature:}
\begin{equation}
    T_E = T_0\!\left(\frac{R_\odot}{2a}\right)^{\!1/2} \approx 278\,\text{K}.
\end{equation}
This is close to the observed average surface temperature; the actual value is
somewhat higher due to the greenhouse effect.
\end{formula}

% ------------------------------------------------------------------
\chapter{Grand Canonical Ensemble}
% ------------------------------------------------------------------

\section{Recap: Gibbs Factor and its Limitations}

When particles are indistinguishable, we correct the classical partition function
by the Gibbs factor $1/N!$, writing $Z_N = Z_1^N/N!$.  This works well in the
\emph{high-temperature limit}, where the typical occupancy of any single quantum
state is much less than one and so two particles are very unlikely to occupy the
same state.  At low temperatures this approximation fails, because occupancies of
order one or more become common, and we must use the correct quantum statistics.

\section{Bosons and Fermions}

Quantum mechanics requires the many-body wavefunction to be \emph{symmetric}
under exchange for \textbf{bosons} (integer spin) and \emph{antisymmetric} for
\textbf{fermions} (half-integer spin).  For two single-particle states $\psi_{k_1}$
and $\psi_{k_2}$:
\begin{align}
    \text{Bosons:}\quad
    \Psi_{k_1k_2}
    &= \frac{1}{\sqrt{2}}\!\left[\psi_{k_1}(1)\psi_{k_2}(2)
       +\psi_{k_2}(1)\psi_{k_1}(2)\right],\\
    \text{Fermions:}\quad
    \Psi_{k_1k_2}
    &= \frac{1}{\sqrt{2}}\!\left[\psi_{k_1}(1)\psi_{k_2}(2)
       -\psi_{k_2}(1)\psi_{k_1}(2)\right].
\end{align}
For fermions the wavefunction vanishes when $k_1=k_2$, so two fermions cannot
occupy the same state.  This is the \textbf{Pauli exclusion principle}.

\begin{fact}
For \textbf{bosons}, the occupation number $n_k$ of any single-particle state
$k$ can be any non-negative integer: $n_k = 0,1,2,\ldots$

For \textbf{fermions}, $n_k\in\{0,1\}$ only.
\end{fact}

\section{General $N$-Particle State}

A general $N$-particle state is labelled by occupation numbers
$|n_1, n_2, \ldots, n_k, \ldots\rangle$, where $n_i$ is the number of particles
in single-particle state $i$ and $\sum_i n_i = N$.

\paragraph{Boson example.}  For $N=3$ spinless (spin-0) bosons in a harmonic
oscillator, the levels are $k=0,1,2,\ldots$ with energies
$E_k=(k+\tfrac{1}{2})\hbar\omega$.  The allowed states are tuples $(n_0,n_1,n_2,\ldots)$
with $\sum n_i = 3$.  There is one state at total energy $\tfrac{3}{2}\hbar\omega$
(all three in $k=0$), one at $\tfrac{5}{2}\hbar\omega$, two at $\tfrac{7}{2}\hbar\omega$
(one particle excited to $k=1$ or $k=2$), and so on.

\begin{center}
[DIAGRAM: Energy-level diagram ($k=0,1,2,3$ horizontal lines) showing four
$N=3$ boson configurations labelled by total occupation tuples; total energies
$\tfrac{3}{2}\hbar\omega$, $\tfrac{5}{2}\hbar\omega$, $\tfrac{7}{2}\hbar\omega$,
$\tfrac{9}{2}\hbar\omega$ indicated below]
\end{center}

\paragraph{Fermion example.}  For $N=3$ spin-$\tfrac{1}{2}$ fermions (each level
2-fold degenerate from spin up/down), the Pauli principle allows at most one
particle per spin-level.  At each total energy there are a countable number of
configurations; the lowest possible energy requires filling the lowest available
spin-states.

\begin{center}
[DIAGRAM: Energy-level diagram for $N=3$ spin-$\tfrac{1}{2}$ fermions, showing
allowed occupation configurations (arrows for spin up/down) at energies
$\tfrac{3}{2}\hbar\omega$ (no state: Pauli blocked), $\tfrac{5}{2}\hbar\omega$
(2 states), $\tfrac{7}{2}\hbar\omega$ (few states), $\tfrac{9}{2}\hbar\omega$
(more states)]
\end{center}

\section{Why the Canonical Partition Function Does Not Factorise}

For a fixed-$N$ system, the canonical partition function is
\begin{equation}
    Z_N = \sum_{\substack{n_1,n_2,\ldots \\ \sum_i n_i = N}}
    e^{-\beta n_1\varepsilon_1}\,e^{-\beta n_2\varepsilon_2}\cdots.
\end{equation}
Because the occupation numbers must satisfy the constraint $\sum_i n_i = N$, they
are not independent, and the sum \emph{cannot} be factored into a product of
independent single-mode sums.  The Gibbs approximation $Z_N \approx Z_1^N/N!$
fails for quantum statistics.  The grand canonical ensemble, which relaxes the
fixed-$N$ constraint, is the natural framework.

\section{The Grand Canonical Ensemble}

\begin{definition}
In the \textbf{grand canonical ensemble} the system is in contact with a reservoir
that fixes both temperature $T$ and chemical potential $\mu$.  The particle number
$N$ is allowed to fluctuate; only the average $\langle N\rangle$ is determined by
$\mu$.  The fixed variables are $V$, $\mu$, $T$.
\end{definition}

The three ensembles are compared as follows.

\begin{center}
[DIAGRAM: Table comparing microcanonical ($V,N,U$ fixed; $T,P,\mu$ derived),
canonical ($V,N,T$ fixed; $U$ fluctuates), and grand canonical ($V,\mu,T$ fixed;
$U$ and $N$ fluctuate)]
\end{center}

\section{New Thermodynamic Potential and Chemical Potential}

Allowing particle exchange modifies the first law:
\begin{equation}
    dU = T\,dS - P\,dV + \mu\,dN.
\end{equation}
The \textbf{chemical potential} $\mu$ is the energy cost of adding one particle
at fixed entropy and volume:
\begin{equation}
    \mu = \pdv{U}{N}\bigg|_{S,V} = -T\pdv{S}{N}\bigg|_{U,V}.
\end{equation}

\paragraph{Physical intuition.}  Consider two subsystems that can exchange both
energy and particles.  Maximising the total entropy with respect to energy exchange
gives $T_1=T_2$ (thermal equilibrium), and maximising with respect to particle
exchange gives
\begin{equation}
    \pdv{S_1}{N_1}\bigg|_{U,V} = \pdv{S_2}{N_2}\bigg|_{U,V}
    \implies \frac{\mu_1}{T_1} = \frac{\mu_2}{T_2},
\end{equation}
so in equilibrium $T$ and $\mu$ are equalised simultaneously.  The chemical
potential plays the same role for particle exchange that temperature plays for
energy exchange.

\begin{center}
[DIAGRAM: Two boxes in thermal and particle contact; arrows showing energy
$U$ and particle number $N$ flow; one box labelled ``system'', the other
``reservoir at $T$, $\mu$'']
\end{center}

Placing the system in contact with a large reservoir at $(T,\mu)$, the system's
energy and particle number both fluctuate.  Let $U_\text{tot}$ and $N_\text{tot}$
denote the conserved totals; then $U_\text{res} = U_\text{tot} - U$ and
$N_\text{res} = N_\text{tot} - N$ are determined by conservation.  The analysis
that follows from expanding the reservoir entropy leads to the grand canonical
probability distribution, which we develop in the next section.
